% ================================================================
%  CHAPTER 5 -- IMPORTANCE SAMPLING: AN AR(1) CASE STUDY
%  Rewrite v4 (discovery-narrative order): bridge from RB -> empirical
%  origin (architecture play, uneven fit) -> the IS idea -> mechanics
%  -> negative result -> diagnosis (why) -> conditions (when) -> contrast.
% ================================================================
\chapter{Importance Sampling}
\label{ch:importance-sampling}

Chapter~\ref{ch:rao-blackwell} reduced the variance of the training
gradient by integrating a conjugate block of the parameter vector out
of the loss. That construction came with a guarantee. The
Rao-Blackwell variance gap~\eqref{eq:variance-gap} is positive
semi-definite, so the construction can never make the gradient noisier,
for any model and any architecture; on Bayesian linear regression it
gave faster convergence and a lower error in every cell of the grid.
This chapter asks whether a second Monte Carlo variance-reduction
technique, \emph{importance sampling}, can help training in the same
way.

The idea came from an experiment rather than from theory. While testing
estimator architectures on a first-order autoregressive model, we
compared the trained network against an exact-posterior reference, and
its accuracy looked uneven across the parameter space. The network
seemed to do worst on the near-non-stationary series at the edges of
the autoregressive coefficient, and on the faint, low-amplitude series
at small noise scale. These are the regions where the realisations are
most distinctive and the process is hardest to simulate, so we suspected
they were also the hardest to estimate. Under a uniform prior, though,
the network sees them no more often than any other part of the
parameter space. Importance sampling offered a way to change that: draw
the training parameters from a proposal that oversamples the difficult
regions, then reweight the loss so that it stays unbiased for the
original prior. If the network's error were concentrated in those
regions, spending more of the simulation budget there should buy
accuracy for the same number of training steps.

The result is negative. Across a family of proposals that oversample
the edges of the parameter space ever more aggressively, the test error
of the trained estimator does not fall; under the most aggressive
proposal it rises, while the effective sample size collapses by more
than an order of magnitude. The negative result is worth reporting
because its cause is specific, and because it contrasts with the
previous chapter. Section~\ref{sec:why} traces the cause to two facts
about the experiment: the trained estimator is already close to the
exact posterior across the whole prior, so little reducible error
remains for any proposal to recover, and the regions we oversampled are
not the ones where its error is largest. Section~\ref{sec:when} sets out
the conditions under which importance sampling does help training, and
why none of them held here. Rao-Blackwellisation never makes training variance worse, for any model and
architecture; importance sampling helps only when the proposal matches
the structure of the problem, and when it does not, the best it can do
is leave the error unchanged.



% ----------------------------------------------------------------
\section{Importance sampling}
\label{sec:is-basics}

Importance sampling is a Monte Carlo technique for estimating an
expectation $\E_{\bm{\theta} \sim p}[f(\bm{\theta})]$ by drawing from a
different distribution $q$ and reweighting:
%
\begin{equation}
  \E_{\bm{\theta} \sim p}[f(\bm{\theta})]
  = \int f(\bm{\theta})\, p(\bm{\theta}) \dd\bm{\theta}
  = \int f(\bm{\theta})\,
      \frac{p(\bm{\theta})}{q(\bm{\theta})}\, q(\bm{\theta}) \dd\bm{\theta}
  = \E_{\bm{\theta} \sim q}\!\left[
      f(\bm{\theta})\,\frac{p(\bm{\theta})}{q(\bm{\theta})}\right].
  \label{eq:is-basic-identity}
\end{equation}
%
The two Monte Carlo estimators,
%
\begin{align}
  \hat{L}_{\mathrm{MC}} &= \frac{1}{N} \sum_{i=1}^{N} f(\bm{\theta}_i), \qquad \bm{\theta}_i \sim p(\bm{\theta}), \\
  \hat{L}_{\mathrm{IS}} &= \frac{1}{N} \sum_{i=1}^{N} w(\bm{\theta}_i)\, f(\bm{\theta}_i), \qquad \bm{\theta}_i \sim q(\bm{\theta}),
\end{align}
%
where $w(\bm{\theta}) = p(\bm{\theta})/q(\bm{\theta})$ is the
\emph{importance weight}, are both unbiased for the same expectation.
The reweighted estimator $\hat{L}_{\mathrm{IS}}$ can have lower variance
than the ordinary estimator $\hat{L}_{\mathrm{MC}}$ when
%
\begin{equation}
  \Var_{\bm{\theta} \sim q}\!\big[w(\bm{\theta})\,f(\bm{\theta})\big]
  < \Var_{\bm{\theta} \sim p}\!\big[f(\bm{\theta})\big].
\end{equation}
%
For a non-negative integrand the proposal that minimises the variance of
$\hat{L}_{\mathrm{IS}}$ is
%
\begin{equation}
  q^{*}(\bm{\theta}) \propto f(\bm{\theta})\,p(\bm{\theta}),
  \label{eq:is-optimal-scalar}
\end{equation}
%
which follows from a Cauchy--Schwarz bound on the second moment of the
reweighted integrand (Appendix~\ref{app:is-optimal-proposal}); under
this proposal the weighted integrand $w(\bm{\theta})f(\bm{\theta})$ is
constant and the variance is zero. Equation~\eqref{eq:is-optimal-scalar}
says to put the proposal mass where the integrand is large, which is
exactly the intuition behind oversampling the hard regions.
Section~\ref{sec:when} returns to this to ask which problems importance
sampling can help; for now it serves as motivation.

In practice the normalising constants of $p$ and $q$ may be unknown, and
the \emph{self-normalised} estimator
%
\begin{equation}
  \hat{r}_{\text{SNIS}}(\bm{\gamma})
  = \frac{\sum_{i=1}^{N} w(\bm{\theta}_i)\,
          L(\bm{\theta}_i, \hat{\bm{\theta}}_{\bm{\gamma}}(\mathbf{y}_i))}
         {\sum_{i=1}^{N} w(\bm{\theta}_i)}
  \label{eq:snis}
\end{equation}
%
is used instead; it is biased at finite $N$ but consistent. To monitor
what a proposal costs we track its \emph{effective sample size}
%
\begin{equation}
  \mathrm{ESS}
  = \frac{\big(\sum_i w(\bm{\theta}_i)\big)^2}{\sum_i w(\bm{\theta}_i)^2}
  = \Big(\textstyle\sum_i \tilde{w}_i^2\Big)^{-1},
  \label{eq:ess}
\end{equation}
%
where $\tilde{w}_i = w(\bm{\theta}_i)/\sum_j w(\bm{\theta}_j)$ are the
normalised weights. It lies between $1$ and $N$ and counts, loosely, how
many of the $N$ reweighted draws carry real weight in the sum. Strictly,
this is the Kish effective sample size of the self-normalised mean: a
convenient proxy for the noise a proposal injects rather than an exact
count of the gradient's effective draws, which
Section~\ref{sec:grad-variance} measures directly. Still, it captures
the price of tilting the proposal away from the prior: when the
effective sample size falls well below $N$, a few draws dominate the
sum, a warning sign that the proposal has been pushed too far from the
prior and may inject more gradient noise than its reweighting removes.

Training a neural Bayes estimator is itself a Monte Carlo integration:
the population loss is the Bayes risk
$r(\bm{\gamma}) = \E_{\bm{\theta}, \mathbf{y}}
[L(\bm{\theta}, \hat{\bm{\theta}}_{\bm{\gamma}}(\mathbf{y}))]$, and each
minibatch estimates its gradient by drawing parameters from the prior.
Drawing $\bm{\theta}$ from a proposal $q$ and reweighting gives the same
risk,
%
\begin{equation}
  r(\bm{\gamma})
  = \int p(\bm{\theta})\,
      \E_{\mathbf{y} \given \bm{\theta}}[L]\dd\bm{\theta}
  = \int \frac{p(\bm{\theta})}{q(\bm{\theta})}\, q(\bm{\theta})\,
      \E_{\mathbf{y} \given \bm{\theta}}[L]\dd\bm{\theta}
  = \E_{\bm{\theta} \sim q}\!\big[w(\bm{\theta})\,
      \E_{\mathbf{y} \given \bm{\theta}}[L]\big].
  \label{eq:is-identity}
\end{equation}
%
The reweighted objective on the right of~\eqref{eq:is-identity} has the
same expectation as the original for every proposal $q$, so the trained
estimator still targets the posterior under the original prior $p$;
reweighting only concentrates the simulation budget where $q$ places
mass. The same identity holds for the gradient that stochastic gradient
descent actually uses, so a proposal changes which parameter values the
gradient is averaged over, not what it estimates.
Section~\ref{sec:when} returns to the question of when this helps.

\section{Application to an AR(1) model}
% ----------------------------------------------------------------
\subsection{The AR(1) model and a recurrent estimator}
\label{sec:ar1-model}

The test model is the first-order autoregressive process of
Example~3, equation~\eqref{eq:ar1}:
$x_t = \rho\,x_{t-1} + \varepsilon_t$ with
$\varepsilon_t \sim \Normal(0, \sigma^2)$, estimand
$\bm{\theta} = (\rho, \sigma)$, priors $\rho \sim \Uniform(-1, 1)$ and
$\sigma \sim \Uniform(0, A)$, and the process initialised from its
stationary distribution $\Normal(0, \sigma^2/(1-\rho^2))$. As
$|\rho| \to 1$ the process approaches non-stationarity, and at
$\rho = 1$ it is a random walk.

Figure~\ref{fig:ar1-examples} shows how the realisations change across
the prior. Persistence grows with the autoregressive coefficient: at
$\rho = 0$ the series is white noise, while near $|\rho| = 1$ it wanders
far from the mean, and the noise scale $\sigma$ sets the amplitude. The
corners of the $(\rho, \sigma)$ box thus produce qualitatively
different, harder-\emph{looking} series than the bulk of the prior. As
the introduction noted, these corners look like where the estimator
should struggle, which makes them the obvious place to spend extra
training draws; Section~\ref{sec:why} shows that for $\rho$ the
appearance is deceptive.

\begin{figure}[t]
  \centering
  \IfFileExists{figures/ar1_examples.png}
    {\includegraphics[width=\textwidth]{ar1_examples}}
    {\fbox{\parbox[c][5cm][c]{\textwidth}{\centering\itshape
      Figure pending: example AR(1) realisations
      (run \texttt{ar1\_examples.py}).}}}
  \caption{Example AR(1) realisations across the prior ($T = 100$, from
  the stationary distribution). Persistence grows with $\rho$ (columns):
  white noise at $\rho = 0$, a near-random-walk at $\rho = 0.99$. The
  noise scale $\sigma$ (rows) sets the amplitude.}
  \label{fig:ar1-examples}
\end{figure}

Because the data~\eqref{eq:ar1} are a time series rather than a set of
exchangeable observations, the permutation-invariant DeepSets
architecture of Section~\ref{sec:deepsets} is inappropriate. The
estimator is instead a gated recurrent unit~\citep{cho2014gru}: a
recurrent encoder processes the series, its hidden states are mean
pooled, and a small multilayer perceptron maps the pooled
representation to $(\hat{\rho}, \hat{\sigma})$, with a $\tanh$ output on
the $\rho$-head and a softplus output on the $\sigma$-head to enforce
the parameter constraints. The estimator is trained under squared-error
loss, so by Proposition~\ref{prop:squared-error} it targets the
posterior mean. The series is fed to the encoder on its raw scale,
without standardisation: the amplitude of an AR(1) realisation is set by
$\sigma$, so rescaling each series to unit variance would erase the
scale that identifies $\sigma$, whereas the raw series lets
the network read $\sigma$ off the amplitude and $\rho$ off the
autocorrelation. Both parameters are then identifiable from the input.

To judge the estimator we need a per-dataset gold standard, and the
AR(1) model supplies one: it is conditionally conjugate in both
parameters, so the Gibbs sampler of Section~\ref{sec:cs-ar1} draws from
the exact posterior, and its post-burn-in average is the posterior-mean
reference against which the trained estimator is judged. Comparing the
two is what makes the loss interpretable: it
separates the error the network could in principle train away from the
error that any estimator on the same data must incur.


% ----------------------------------------------------------------
\subsection{A family of proposals}
\label{sec:is-proposals}

With the goal of training our neural network to better fit the edges of the prior, we need a family of proposals that oversample those edges and a single knob that controls how aggressively. For a parameter with bounded support the Beta
distribution provides exactly this. The $\mathrm{Beta}(\alpha, \alpha)$
density on $[0,1]$ is symmetric, and the single parameter $\alpha$
controls where it places mass: $\alpha = 1$ is uniform, $\alpha < 1$ is
U-shaped and concentrates mass at the boundaries, and $\alpha > 1$ is
bell-shaped and concentrates it in the centre. The proposals that target
the difficult-looking boundary take $\alpha < 1$, with $\alpha = 1$ as
the uniform baseline. We also include two proposals on the opposite tilt,
at $\alpha = 1.5$ and $\alpha = 2$; what they do to each coordinate is set
out below.

For the autoregressive coefficient, a draw
$\tilde{\rho} \sim \mathrm{Beta}(\alpha, \alpha)$ on $(0,1)$ is mapped to
$(-1,1)$ by $\rho = 2\tilde{\rho} - 1$; by the change-of-variables
formula the proposal density and importance weight are
%
\begin{equation}
  q(\rho) = \tfrac{1}{2}\,
    \mathrm{Beta}\!\Big(\tfrac{\rho+1}{2};\, \alpha, \alpha\Big),
  \qquad
  w(\rho) = \frac{p(\rho)}{q(\rho)}
  = \frac{1}{\mathrm{Beta}\!\big(\tfrac{\rho+1}{2};\, \alpha, \alpha\big)}.
  \label{eq:is-weight}
\end{equation}
%
At $\alpha = 1$ the Beta density is identically one and $w \equiv 1$,
recovering uniform training. The noise scale is treated differently.
Small $\sigma$ gives a faint, near-white-noise series that looks hard to
pin down, so we oversample only that low-noise tail, as a one-sided test
of whether concentrating draws on a plausibly hard region buys accuracy
there. An asymmetric proposal
$\tilde{\sigma} \sim \mathrm{Beta}(\alpha, 1)$, mapped by
$\sigma = A\tilde{\sigma}$, does this:
%
\begin{equation}
  q(\sigma) = \tfrac{1}{A}\,
    \mathrm{Beta}\!\Big(\tfrac{\sigma}{A};\, \alpha, 1\Big),
  \qquad
  w(\sigma) = \frac{1}{\mathrm{Beta}\!\big(\tfrac{\sigma}{A};\, \alpha, 1\big)},
  \label{eq:is-weight-sigma}
\end{equation}
%
and the joint weight is the product of the two factors. Run at
$\alpha > 1$ instead, the two proposals reverse: the $\rho$-proposal
concentrates on the interior near $\rho = 0$ and the $\sigma$-proposal on
large $\sigma$. We include two such cases, at $\alpha = 1.5$ and
$\alpha = 2$, as the opposite tilt, to
check whether aiming at those regions does any better than aiming at the
boundary. Smaller $\alpha$
oversamples the targeted regions more aggressively but makes $w$ more
variable, which lowers the effective sample size through~\eqref{eq:ess}.
That tension, between concentrating the draws and keeping enough of them
informative, is what the sweep below probes.
Figure~\ref{fig:is-proposals} shows the mechanism: as $\alpha$ falls the
$\rho$-proposal piles mass at both ends $|\rho| \to 1$ and the
$\sigma$-proposal piles mass at small $\sigma$, while the compensating
weights spread over more orders of magnitude.
Figure~\ref{fig:is-samples} makes the same reshaping concrete in the
$(\rho, \sigma)$ plane.

\begin{figure}[t]
  \centering
  \IfFileExists{figures/proposals.png}
    {\includegraphics[width=\textwidth]{proposals}}
    {\fbox{\parbox[c][5cm][c]{\textwidth}{\centering\itshape
      Figure pending: proposal densities and importance weights
      (run \texttt{plot\_proposals.py}).}}}
  \caption{The Beta proposal family. \textbf{(a)} The $\rho$-proposal
  $\mathrm{Beta}(\alpha,\alpha)$: $\alpha < 1$ oversamples the boundaries
  $|\rho| \to 1$, $\alpha > 1$ the interior, $\alpha = 1$ is uniform.
  \textbf{(b)} The $\sigma$-proposal $\mathrm{Beta}(\alpha,1)$:
  $\alpha < 1$ oversamples small $\sigma$, $\alpha > 1$ large $\sigma$.
  \textbf{(c, d)} The corresponding importance weights $w = p/q$.}
  \label{fig:is-proposals}
\end{figure}

\begin{figure}[t]
  \centering
  \IfFileExists{figures/samples_weighted.png}
    {\includegraphics[width=\textwidth]{samples_weighted}}
    {\fbox{\parbox[c][5cm][c]{\textwidth}{\centering\itshape
      Figure pending: $(\rho,\sigma)$ draws coloured by weight
      (run \texttt{plot\_proposals.py}).}}}
  \caption{Training draws in the $(\rho,\sigma)$ plane, coloured by
  importance weight $w$ (log scale). \textbf{(a)} The uniform prior
  ($\alpha = 1$): even draws, unit weights. \textbf{(b)} An aggressive
  proposal ($\alpha = 0.2$): draws cluster on the $|\rho| \to 1$ edges and
  the small-$\sigma$ floor (small weights, dark), leaving the interior and
  large $\sigma$ sparse (large weights, bright).}
  \label{fig:is-samples}
\end{figure}


% ----------------------------------------------------------------
\subsection{Results}
\label{sec:is-results}

Estimators were trained at six proposal strengths,
$\alpha \in \{2.0, 1.5, 1.0, 0.5, 0.3, 0.1\}$: the uniform baseline at
$\alpha = 1$, three increasingly aggressive boundary proposals at
$\alpha < 1$, and two proposals on the opposite tilt at $\alpha = 1.5$ and
$\alpha = 2$. The test error was
broken down over the whole prior and over three regions of interest: the
near-non-stationary boundary $|\rho| > 0.8$, the low-noise region
$\sigma < 0.25A$, and the high-noise region $\sigma > 0.75A$.
Table~\ref{tab:is-sweep} reports the mean squared error in each region
together with the mean effective sample size at each $\alpha$.

\begin{table}[t]
  \centering
  \caption{Test MSE by proposal strength $\alpha$ and region, with the
  mean effective sample size (batch size $256$). Regions: the
  near-non-stationary boundary $|\rho| > 0.8$, low noise
  $\sigma < 0.25A$, and high noise $\sigma > 0.75A$.}
  \label{tab:is-sweep}
  \begin{tabular}{c rrrr r}
    \toprule
    $\alpha$ & all & $|\rho|>0.8$ & low $\sigma$ & high $\sigma$ & ESS \\
    \midrule
    $2.0$ & $0.0181$ & $0.0196$ & $0.0161$ & $0.0245$ & $\phantom{0}93.1$ \\
    $1.5$ & $0.0152$ & $0.0123$ & $0.0121$ & $0.0198$ & $172.7$ \\
    $1.0$ & $0.0136$ & $0.0100$ & $0.0087$ & $0.0187$ & $256.0$ \\
    $0.5$ & $0.0141$ & $0.0102$ & $0.0099$ & $0.0188$ & $155.9$ \\
    $0.3$ & $0.0139$ & $0.0100$ & $0.0084$ & $0.0181$ & $\phantom{0}78.8$ \\
    $0.1$ & $0.0163$ & $0.0128$ & $0.0112$ & $0.0217$ & $\phantom{0}13.4$ \\
    \bottomrule
  \end{tabular}
\end{table}

The result is negative. Across the whole prior the test error is
essentially flat for the baseline and the two milder boundary proposals
($0.0136$, $0.0141$, $0.0139$) and rises at both extremes: climbing through
$\alpha = 1.5$ ($0.0152$) to $0.0181$ at $\alpha = 2$ on one side, and to
$0.0163$ at $\alpha = 0.1$ on the other. No proposal improves on the
uniform baseline. The
regional breakdown (Figure~\ref{fig:is-region}) shows why there is
nothing to redistribute. The low-$\sigma$ region the boundary proposals
oversample is already the easiest at baseline ($0.0087$, the smallest of
the four), and flooding it with draws leaves it there; the high-$\sigma$
region they starve is the hardest ($0.0187$) and stays so; and the
$|\rho| > 0.8$ boundary the $\rho$-proposal targets is among the regions
where the $\rho$-error is smallest. The interior tilts at $\alpha = 1.5$
and $\alpha = 2$ aim the other way, towards the $\rho$-interior and high
$\sigma$, and fare no better: their error is higher in every region, the
high-$\sigma$ region they now oversample included ($0.0245$ at
$\alpha = 2$). The estimator already handles every
region about as well as it can, so reweighting finds no slack to recover,
whichever way it aims. At $\alpha = 0.1$ this is starkest: every region
degrades together, the flooded low-$\sigma$ tail included (it rises from
$0.0087$ to $0.0112$), as the collapsed effective sample size leaves the
gradient too noisy to train well.

\begin{figure}[t]
  \centering
  \IfFileExists{figures/region_breakdown.png}
    {\includegraphics[width=0.85\textwidth]{region_breakdown}}
    {\fbox{\parbox[c][6cm][c]{0.85\textwidth}{\centering\itshape
      Figure pending: per-region MSE bars
      (re-export \texttt{region\_breakdown} from the AR(1) sweep).}}}
  \caption{Test MSE by region for each proposal strength $\alpha$. The
  baseline and the two milder boundary proposals ($\alpha = 1, 0.5, 0.3$)
  sit lowest and together; the interior tilt $\alpha = 1.5$ is intermediate,
  and the extremes $\alpha = 0.1$ and $\alpha = 2$ are worse in every
  region.}
  \label{fig:is-region}
\end{figure}

What does fall is the effective sample size: it peaks at $256$ for the
uniform baseline and drops away as the proposal tilts in either
direction, to $173$ at $\alpha = 1.5$ and $93$ at $\alpha = 2$ on one side
and $13$ at $\alpha = 0.1$ on the other, as the
increasingly variable weights~\eqref{eq:is-weight} let a handful of draws
dominate the reweighted sum. This is the proposal's one reliable effect,
and it is a cost, not a benefit. The training curves
(Figure~\ref{fig:is-curves}) show its consequence: the three milder
proposals converge together near the floor, while $\alpha = 1.5$ settles a
little above them and the two extremes, $\alpha = 0.1$ and $\alpha = 2$,
sit higher still throughout, jitterier and settling at a
higher loss rather than the same one.

\begin{figure}[t]
  \centering
  \IfFileExists{figures/training_curves.png}
    {\includegraphics[width=0.8\textwidth]{training_curves}}
    {\fbox{\parbox[c][6cm][c]{0.8\textwidth}{\centering\itshape
      Figure pending: validation MSE training curves
      (re-export \texttt{training\_curves} from the AR(1) sweep).}}}
  \caption{Seed-averaged validation MSE over training for the six
  proposals. The three milder settings ($\alpha = 1, 0.5, 0.3$) converge
  together near the floor; the interior tilt $\alpha = 1.5$ sits a little
  above them, and the extremes $\alpha = 0.1$ and $\alpha = 2$ stay higher
  still, jitterier and settling higher.}
  \label{fig:is-curves}
\end{figure}
\subsection{The gradient variance, measured directly}
\label{sec:grad-variance}

The previous section read the negative result off two indirect quantities: the
test error of the trained estimator, which no proposal lowered, and the
effective sample size, which every proposal lowered. Neither is the quantity
importance sampling is meant to reduce. What training actually consumes is the
gradient: by~\eqref{eq:is-identity} each minibatch estimates
$\nabla_{\bm{\gamma}} r$ by reweighting draws from the proposal, and the
speed of stochastic gradient descent is governed by the variance of that
estimate, not by the variance of the loss nor by the effective sample size on its
own. This section measures the gradient variance directly, and so asks of each
proposal the question the chapter has thus far answered only by proxy: did
reweighting make the gradient less noisy, or more?

The measurement holds the estimator fixed. At a frozen set of weights
$\bm{\gamma}$ and a proposal $q_{\alpha}$, a single minibatch of the common size
$B = 256$ produces the self-normalised gradient
%
\begin{equation}
  \mathbf{g}
  = \nabla_{\bm{\gamma}} \sum_{i=1}^{B} \tilde{w}_i\,
      L(\bm{\theta}_i, \hat{\bm{\theta}}_{\bm{\gamma}}(\mathbf{y}_i)),
  \qquad \bm{\theta}_i \sim q_{\alpha},
  \label{eq:grad-est}
\end{equation}
%
exactly the estimator the training loop applies in~\eqref{eq:snis}. Drawing $M$
independent minibatches, each of the same size $B$ so that every $\alpha$ spends
the identical simulation budget per step, the total gradient variance is
estimated by
%
\begin{equation}
  V(\alpha)
  = \Tr\bigl(\Cov(\mathbf{g})\bigr)
  = \frac{1}{M} \sum_{m=1}^{M}
      \big\lVert \mathbf{g}^{(m)} - \bar{\mathbf{g}} \big\rVert_2^2,
  \qquad
  \bar{\mathbf{g}} = \frac{1}{M} \sum_{m=1}^{M} \mathbf{g}^{(m)},
  \label{eq:grad-var}
\end{equation}
%
the spread of the minibatch gradient about its own mean, summed over every weight
of the network. The quantity reported is the ratio $V(\alpha)/V(1)$ to the
uniform baseline: a value below one is a genuine variance reduction, a value
above one is added noise, and the baseline sits at one by construction. Because
the variance depends on the weights at which it is taken, the comparison is made
at the same network for every proposal: a fresh random initialisation, fixed and
identical across the sweep. Each ratio is averaged over three measurement seeds.

The result is unambiguous (Figure~\ref{fig:grad-variance}). The gradient variance
is smallest at the uniform baseline and rises under every proposal; no setting of
$\alpha$ reduces it, and the inflation grows monotonically as the proposal tilts
further from uniform in either direction. The milder proposals raise it only
gently --- by a factor of about $1.3$ at $\alpha = 0.5$, $2.2$ at $\alpha = 0.3$,
and $2.3$ at the opposite tilt $\alpha = 1.5$ --- while the extremes raise it
sharply, by roughly $14\times$ at $\alpha = 0.1$ and $8.6\times$ at $\alpha = 2$.
This is the direct counterpart of the test-error and effective-sample-size
findings of Section~\ref{sec:is-results}: the one reliable effect of reweighting
here is to make the gradient noisier. No proposal in the Beta family reduced the
variance that training consumes; Section~\ref{sec:why} explains why none
could.
%
\begin{figure}[t]
  \centering
  \IfFileExists{figures/grad_variance.png}
    {\includegraphics[width=0.8\textwidth]{grad_variance}}
    {\fbox{\parbox[c][6cm][c]{0.8\textwidth}{\centering\itshape
      Figure pending: gradient-variance ratio $V(\alpha)/V(1)$ vs $\alpha$
      (run \texttt{grad\_variance.py}).}}}
  \caption{Per-minibatch gradient variance by proposal strength, as the ratio
  $V(\alpha)/V(1)$ to the uniform baseline, measured at a fresh random
  initialisation and averaged over three seeds (error bars show the seed spread).
  The dashed line marks the uniform baseline at one; every proposal lies above
  it, so none reduces the variance, and the two extremes ($\alpha = 0.1$ and
  $\alpha = 2$) inflate it by roughly an order of magnitude.}
  \label{fig:grad-variance}
\end{figure}


% ----------------------------------------------------------------
\subsection{Why reweighting cannot help here}
\label{sec:why}

Reweighting does not change what the training
objective is, only how noisily each minibatch estimates
it~\eqref{eq:is-identity}, so it can speed training only by making the
gradient less noisy. For that to buy anything, two things must hold: the
estimator must have error left to remove, and the proposal must put its
draws where that error is. On the AR(1) problem neither holds, and both
facts come from comparing the trained network with the exact posterior.

The estimator is already about as accurate as any estimator could be.
Table~\ref{tab:gru-vs-gibbs} compares the GRU with the Gibbs posterior
mean, the Bayes-optimal estimator for this model, on the same test sets.
The GRU's error exceeds the Gibbs error by $22\%$ on the autoregressive
coefficient and $13\%$ on the noise scale.

\begin{table}[t]
  \centering
  \caption{Test MSE of the GRU estimator and the Gibbs posterior-mean
  (Bayes) reference, by parameter, on a held-out test set.}
  \label{tab:gru-vs-gibbs}
  \begin{tabular}{l rr}
    \toprule
             & MSE in $\rho$ & MSE in $\sigma$ \\
    \midrule
    Gibbs posterior mean (Bayes) & $0.0073$ & $0.0056$ \\
    GRU estimator                & $0.0089$ & $0.0063$ \\
    \bottomrule
  \end{tabular}
\end{table}

The error on any one dataset has two parts. One is the posterior
uncertainty that no estimator can remove: even the Bayes-optimal
posterior mean misses, because the data do not pin the parameters down
exactly. The other is the excess the network carries on top, the gap
between it and that optimum. Reweighting cannot touch the first part, and
it recovers the second only where training is held back by gradient
noise, the one thing a cleaner gradient improves. That excess is small in
absolute terms (Table~\ref{tab:gru-vs-gibbs}), and
Figure~\ref{fig:sigma-failure} shows what it is: on the hardest series the
GRU and the Gibbs sampler miss together, so the error there is posterior
uncertainty, not a deficit a better-aimed gradient could close. Little of
the error anywhere is the kind a proposal could recover.

\begin{figure}[t]
  \centering
  \IfFileExists{figures/gibbs_vs_nn.png}
    {\includegraphics[width=\textwidth]{gibbs_vs_nn}}
    {\fbox{\parbox[c][5.5cm][c]{\textwidth}{\centering\itshape
      Figure pending: GRU versus Gibbs estimates by parameter
      (run \texttt{gibbs\_vs\_nn.py}).}}}
  \caption{Estimates from the baseline GRU (blue) and the Gibbs posterior
  mean (grey) against the truth, for $\rho$ (left) and $\sigma$ (centre);
  the worst-MSE cases are red. Right: one of them as a realisation, an
  extreme series with $\rho \approx -0.99$ and $\sigma$ at the top of its
  range. The GRU recovers $\rho$ exactly, though the series looks the
  hardest of all; both it and Gibbs underestimate the extreme $\sigma$, so
  that residual is shared with the Bayes-optimal sampler, not something the
  network could train away.}
  \label{fig:sigma-failure}
\end{figure}

The little error that remains sits in the opposite place from where the
proposals aim. Figure~\ref{fig:binned-mse}
bins the test error by each parameter. The error in $\rho$ is largest in
the interior, near $\rho = 0$, and smallest at the boundaries
$|\rho| \to 1$, the very edge the $\rho$-proposal floods; it is also
raised where $\sigma$ is small. The error in $\sigma$ runs the other way,
growing with $\sigma$, so the high-noise end carries most of it while the
low-$\sigma$ tail the $\sigma$-proposal floods is the cleanest part.

There is a clean reason the boundary is the easy end for $\rho$.
Estimating $\rho$ is a regression of each $x_t$ on its lag $x_{t-1}$, so
the posterior variance of the slope is the noise variance over the spread
of the predictor, $\Var(\rho \mid \sigma^2, \mathbf{x}) = \sigma^2/S$ with
$S = \sum_t x_{t-1}^2$ (Section~\ref{sec:cs-ar1}). The series starts from
its stationary distribution, of variance $\sigma^2/(1-\rho^2)$, so
$S \approx T\sigma^2/(1-\rho^2)$ and the posterior variance is
$\approx (1-\rho^2)/T$. As $|\rho| \to 1$ the series makes larger, more
persistent excursions, and this variance vanishes: the data pin $\rho$
down most tightly exactly where the series looks most dramatic. The
motivating intuition mistook one kind of difficulty for another: a
near-non-stationary series has large
swings that carry the most information about $\rho$, so it is the
\emph{easiest} to estimate $\rho$ from. The hard-to-simulate end and the
hard-to-estimate end are opposite.

\begin{figure}[t]
  \centering
  \IfFileExists{figures/binned_mse.png}
    {\includegraphics[width=\textwidth]{binned_mse}}
    {\fbox{\parbox[c][5.5cm][c]{\textwidth}{\centering\itshape
      Figure pending: test MSE binned by parameter
      (run \texttt{extra\_figures.py}).}}}
  \caption{Test error binned by $\rho$ (top row) and by $\sigma$ (bottom
  row), for the baseline (blue), the boundary proposal $\alpha = 0.3$
  (red), and the opposite tilt $\alpha = 2$ (green). The error lives in
  the $\rho$-interior and at large $\sigma$, not at the boundary or
  low-$\sigma$ tail the proposals target; neither overlay improves on the
  baseline.}
  \label{fig:binned-mse}
\end{figure}

The $\alpha = 2$
proposal puts its draws in the $\rho$-interior and at high $\sigma$, where
the error genuinely is largest, and it is no better than the rest
(Table~\ref{tab:is-sweep}), because the error there is the irreducible
kind the Gibbs sampler shares. So no setting of the Beta family helps:
where the error is large it cannot be removed, and where it can be removed
it is already small.

Against these absent gains stands a cost that is always paid. The more a
proposal concentrates its draws, the more variable its
weights~\eqref{eq:is-weight}, the lower the effective sample
size~\eqref{eq:ess}, and the noisier the gradient left by the
self-normalised objective~\eqref{eq:snis}. With no gain to offset it, this
cost is the whole effect: at $\alpha = 0.1$ it shows up as the
across-the-board degradation of Section~\ref{sec:is-results}. Reweighting
here can only add noise, and the harder it is pushed, the more it adds.


% ----------------------------------------------------------------
\section{When importance sampling can help}
\label{sec:when}


The negative result is therefore a verdict on a mismatch, not on importance sampling in general. Three conditions separate the problems in which importance sampling can help from those in which it cannot.

The first is that there must be reducible error left to recover. Section~\ref{sec:why} showed that little such error remains in the AR(1) experiment: the trained GRU is already close to the Gibbs posterior-mean reference. The remaining test error is therefore dominated by posterior uncertainty rather than by a failure of optimisation. Reweighting the training distribution cannot remove this irreducible part of the Bayes risk.

The second condition is that the proposal must target where the useful gradient signal is concentrated, not merely where the raw loss is large. For squared-error loss,
\[
L = \|\hat{\bm{\theta}}_{\bm{\gamma}}(\mathbf{y})-\bm{\theta}\|_2^2,
\qquad
\nabla_{\bm{\gamma}} L = 2\left(\nabla_{\bm{\gamma}}\hat{\bm{\theta}}_{\bm{\gamma}}(\mathbf{y})\right)^{\!\top}\left(\hat{\bm{\theta}}_{\bm{\gamma}}(\mathbf{y})-\bm{\theta}\right).
\]

Thus, ignoring variation in the estimator sensitivity $(\|\nabla_{\bm{\gamma}}\hat{\bm{\theta}}_{\bm{\gamma}}(\mathbf{y})\|)$, the gradient magnitude scales like $\sqrt{L}$, not like $L$ itself. A proposal designed to reduce gradient variance should therefore be guided by the size of the gradient contribution, roughly $q \propto p\sqrt{L}$, rather than by the loss alone. More importantly, the relevant loss is the reducible excess over the Bayes estimator. Large loss caused by posterior uncertainty produces noisy residuals, but not a systematic gradient direction that the network can train away.

The third condition is that the variance affected by the proposal must be the dominant source of training noise. Importance sampling changes how often different parameter values $\bm{\theta}$ are drawn, but it does not remove the randomness from simulating a fresh dataset $(\mathbf{y}\mid\bm{\theta})$. If that conditional simulation noise dominates the minibatch gradient, then even a well-aimed proposal can only produce a limited gain.

The proposals used here were fixed in advance by the single parameter
$\alpha$, which cannot know where the recoverable gradient signal lies; on
this problem it aimed badly. One direction for future work is to adapt the
proposal during training, estimating where the gradient is largest as
training proceeds and reshaping $q \propto p\sqrt{L}$ towards it, so that
the simulation budget follows the signal rather than a fixed guess.
Whether such a scheme can realise even the bounded gains these conditions
allow, without its own overhead and instability, is an open question.
A second is to test the method on a problem built to meet the conditions rather than to fail them, where the loss is sharply uneven, the estimator's sensitivity varies in the same direction, and the
residual error is genuinely reducible. That would be the fair positive
counterpart to the AR(1) case study, and the place to further investigate whether
importance sampling can meaningfully help neural Bayes training.


% ----------------------------------------------------------------
\section{Discussion}
\label{sec:is-discussion}

The experiment returns a negative result, and the negative result is the
point. On the AR(1) model, importance sampling with Beta proposals leaves
the trained estimator's test error unchanged for the milder proposals and
raises it at both extremes, all while collapsing the effective sample size
by more than an order of magnitude. Section~\ref{sec:why} explains why:
the estimator is already close to the Bayes reference everywhere, so
little reducible error remains for any proposal to recover, and the
regions the boundary proposals target are the ones the estimator already
handles best. The motivating intuition confused the conspicuous regions,
the ones that are dramatic to look at and awkward to simulate, with the
difficult ones; for the autoregressive coefficient the two are opposite.

Section~\ref{sec:when} turns that failure into a specification: importance
sampling helps training only when there is reducible error to recover,
when the proposal is aimed where the gradient signal actually
concentrates, and when the simulation noise is not what limits training. A
scale-estimation problem can be built to meet these conditions, where the
gain is real if modest; the AR(1) study met none of them where its
proposals aimed.

This is the genuine contrast with Chapter~\ref{ch:rao-blackwell}. The
Rao-Blackwell variance gap~\eqref{eq:variance-gap} is positive
semi-definite for every model and every architecture: the technique
cannot make training variance worse, and on Bayesian linear regression it made
training reliably and substantially better. Importance sampling carries
no such guarantee. It helps only when the proposal is well matched and
gradient noise is the binding constraint; when the difficulty lies
elsewhere, as here, the best it can do is leave the error untouched, and
pushed harder it makes the error worse, while its cost in effective sample
size is paid either way. Both are instances of Monte Carlo variance
reduction applied to estimator training, but they differ in reliability:
one unconditional, the other contingent on a match between the proposal
and the structure of the problem. That difference is the main lesson of
this chapter.